\scanpage{54}
% archive foliation: 53
2. \uline{Courbes relatives de type} $(0,4)$.

Soit $X$ un fibré en droites projectives sur $S$, et
\[
(9)\qquad\underline I\subset X
\]
un diviseur relatif, étale de degré $4$. Considérons
\[
(10)\qquad\underline{\mathfrak S}_{\underline I}=\underline{\operatorname{Aut}}_S(I)
\]
qui est un schéma en groupes fini étale sur $S$, loc isom. à $(\mathfrak S_4)_S$. Il lui correspond un unique sous-groupe \inserted{distingué} $\underline V_{\underline I}$, loc. isom. à $(\mathbb D_2)_S$, \deleted{de telle faç\unsure{on}} Mais $\underline I$ est un torseur sous $\underline V_I$.

\uline{Proposition}\quad L'opération de $\underline V_{\underline I}$ sur $\underline I$ se prolonge (de façon évidemment unique) en une opération de $\underline V_{\underline I}$ sur $X$.

\uline{NB} La donnée de $\underline I$, fini étale de degré $4$ sur $S$, équivaut à celle de $\underline V=\underline V_{\underline I}$ sur $S$ (ou encore de $\underline V_{\underline I}^*=\underline J_{\underline I}$, fini étale de degré $3$ sur $S$), plus la donnée d'un torseur $\underline I$ sous $\underline V^*$. Plonger alors $\underline I$ dans un fibré en droites projectives $X$, revient à faire opérer fidèlement $\underline V$ sur $X$ (chose étudiée en § 1), et de plus, de \unsure{se}
